\documentclass{article}
\usepackage[T1]{fontenc}
\usepackage{lmodern}
\usepackage{graphicx}
\usepackage{amsmath,amssymb}
\usepackage[a4paper,margin=2cm]{geometry}
\usepackage[absolute,overlay]{textpos}
\setlength{\TPHorizModule}{1mm}
\setlength{\TPVertModule}{1mm}
\usepackage[indonesian]{babel}
\usepackage{hyperref}
\setlength{\emergencystretch}{2em}
% unit: Halaman 64

\begin{document}

% === Halaman 64 (python/native) ===

64

• A5 terdiri dari 3 buah LFSR , masing-masing panjangnya 19, 22, 

dan 23 bit (total = 19 + 22 + 23 = 64). 

• Bit-bit di dalam register diindeks dimana bit paling tidak penting 

(LSB) diindeks dengan 0 (elemen paling kanan). 

• Luaran (output) dari A5 adalah hasil XOR dari ketiga buah LFSR ini. 

• A5 menggunakan tiga buah kendali detak (clock) yang variabel:  

  - Bit ke-8 pada register 1. Bit-bit detak pada bit 13, 16, 17, dan 18

 - Bit ke-10 pada register 2. Bit-bit detaknya adalah pada bit 20 dan 21

 - Bit ke-10 pada register 3. Bit-bit detaknya adalah pada bit 7, 20, 21, dan 22

\end{document}
